\documentclass[10pt, twocolumn, landscape, a4paper]{article}

\usepackage[margin=1in]{geometry}
\usepackage[utf8]{inputenc}
\usepackage[spanish]{babel}
\usepackage{amssymb, amsmath, amsbsy}
\usepackage{verbatim}
\usepackage{mathpazo}
\usepackage{tasks}
\usepackage{enumerate}
\usepackage{graphicx}
\usepackage[pdftex]{hyperref}
\hypersetup{pdftitle={valor numérico de polinomios},pdfauthor={Luis Carrillo Gutiérrez}}

\renewcommand{\familydefault}{\sfdefault}

\begin{document}
\subsubsection*{Valor numérico de polinomios}
% El valor numérico de una expresión algebraica es el resultado de sustituir las letras o variables por números y operar después. Una expresión algebraica puede tener tantos valores numéricos como valores diferentes les demos a las letras.
\begin{itemize}
\item{Calcular el valor numérico del polinomio : \quad$P(x) = 3x^2 - 5x + 1$ \\ para $x = -2$\quad y \quad$x = \dfrac1{4}$}
\item{Calcular el valor numérico de la expresión algebraica : $Q(x,\,y) = a^2 + 2ab - b$ en los casos : 
\begin{itemize}
\item{$a = -3,\qquad b = 4$}
\item{$a = \dfrac1{3},\qquad b = -\dfrac1{2}$}
\end{itemize}
}
\item{Calcular el valor numérico de los siguientes polinomios.
\begin{itemize}
\item{$R(x) = 5x^3 - 3x^2 + 6x$ \qquad para $x = -1$}
\end{itemize}
}
\end{itemize}
\end{document}